\chapter{Persamaan Diferensial Biasa}\label{ch:b3:ode}

Tahun ke-2 menyelesaikan persamaan diferensial linear dan menyatakan
teorema Cauchy--Lipschitz; sedangkan bab ini membuktikannya --- dua kali:
keberadaan dan ketunggalannya lewat titik tetap Banach, serta
struktur globalnya lewat teori \omterm{thm:b3:ode:maximal}{solusi maksimal} dan
teorema \omterm{thm:b3:ode:maximal}{pelolosan dari kompak}. Lalu teori linearnya dibangun
kembali di atas landasan yang jujur (\omterm{thm:b3:ode:linearstructure}{resolven}, \omterm{thm:b3:ode:linearstructure}{Wronskian}, \omterm{thm:b3:ode:matrixexp}{eksponensial
matriks}, Duhamel), dan paruh kedua bab ini membuka
teori \emph{kualitatifnya} --- yakni \omterm{ex:b3:ode:planeclassification}{aliran}, kesetimbangan,
\omterm{thm:b3:ode:lyapunov}{fungsi Lyapunov}, dan \omterm{thm:b3:ode:linearization}{kestabilan lewat pelinearan}: bagaimana
memahami solusi yang tak akan pernah kita hitung. Sedangkan bandulnya, pada
soal akhir pekan, adalah studi kasus yang abadi. Sepanjang bab ini,
$U \subseteq \R\times\R^d$ terbuka dan $f \colon U \to \R^d$
\omterm{def:b3:topology:continuity}{kontinu}; dan sebuah \emph{solusi} $x' = f(t, x)$ adalah
pemetaan $\mathcal C^1$ $x \colon I \to \R^d$ (dengan $I$ sebuah selang)
yang grafiknya di $U$ dan memenuhi persamaannya.

\section{Cauchy--Lipschitz}

\begin{definition}\label{def:b3:ode:lipschitz}
Fungsi $f$ disebut \emph{Lipschitz lokal terhadap $x$}\index{Lipschitz lokal}
jika setiap titik $U$ berpersekitaran $V$ dan berkonstanta
$L$ dengan $\norm{f(t, x_1) - f(t, x_2)} \leq L\norm{x_1 -
x_2}$ untuk $(t, x_1), (t, x_2) \in V$. Jika $f$ bersifat $\mathcal
C^1$ (atau sekadar $\partial_xf$ ada dan \omterm{def:b3:topology:continuity}{kontinu}), maka ia
Lipschitz lokal terhadap $x$: sebab pada sebuah \omterm{def:b3:topology:topology}{persekitaran} cembung
\omterm{def:b3:topology:compact}{kompak}, lewat ketaksamaan nilai rata-rata dengan $L =
\sup\vertiii{\partial_xf}$.
\end{definition}

\begin{theorem}[Cauchy--Lipschitz, lokal]
\label{thm:b3:ode:cauchylipschitz}
Misalkan $f$ \omterm{def:b3:topology:continuity}{kontinu} dan \omterm{def:b3:ode:lipschitz}{Lipschitz lokal} terhadap $x$, dan $(t_0,
x_0) \in U$. Maka ada $T > 0$ sedemikian sehingga soal Cauchy
\[
x' = f(t, x), \qquad x(t_0) = x_0
\]
mempunyai tepat satu solusi pada $[t_0 - T, t_0 +
T]$.\index{teorema Cauchy--Lipschitz}
\end{theorem}

\begin{proof}
Pilihlah $a, b > 0$ dengan $Q = [t_0 - a, t_0 + a]\times\bar
B(x_0, b) \subseteq U$, yang $\norm f \leq M$ padanya dan $f$ bersifat
Lipschitz-$L$ terhadap $x$. Sebuah fungsi $\mathcal C^1$ merupakan solusi
jika dan hanya jika ia memenuhi persamaan integral
\[
x(t) = x_0 + \int_{t_0}^{t}f\bigl(s, x(s)\bigr)\,\dd s
\]
(lewat teorema fundamental kalkulus, dalam kedua arahnya). Lalu misalkan $T =
\min\bigl(a, \frac bM, \frac1{2L}\bigr)$, $I = [t_0 - T, t_0
+ T]$, dan
\[
\mathcal E = \{x \in \mathcal C(I, \R^d) : \norm{x(t) - x_0}
\leq b\ \text{pada } I\},
\]
yakni sebuah himpunan bagian tertutup ruang Banach $(\mathcal C(I, \R^d),
\norm\cdot_\infty)$: jadi \omterm{def:b3:complete:complete}{lengkap}
(\cref{def:b3:complete:complete}). Definisikan $\Phi(x)(t) = x_0 +
\int_{t_0}^tf(s, x(s))\dd s$: maka untuk $x \in \mathcal E$,
$\norm{\Phi(x)(t) - x_0} \leq M\abs{t - t_0} \leq MT \leq b$
--- sehingga $\Phi$ memetakan $\mathcal E$ ke dirinya --- dan untuk $x, y \in
\mathcal E$:
\[
\norm{\Phi(x)(t) - \Phi(y)(t)} \leq \Bigl|\int_{t_0}^t
L\,\norm{x(s) - y(s)}\,\dd s\Bigr| \leq LT\,\norm{x -
y}_\infty \leq \tfrac12\norm{x - y}_\infty :
\]
yakni sebuah kontraksi. Lalu titik tetap Banach
(\cref{thm:b3:complete:banach}) memberikan satu-satunya titik tetap di
$\mathcal E$: yakni keberadaannya, beserta ketunggalannya di antara solusi
yang tinggal di $\bar B(x_0, b)$ --- tetapi setiap solusi pada $I$ tinggal
di sana (sebab $\norm{x(t) - x_0} \leq M\abs{t - t_0} \leq b$ selama
grafiknya tetap di $Q$, lewat argumen \omterm{def:b3:topology:continuity}{kekontinuannya}):
sehingga ketunggalannya pada $I$.
\end{proof}

\begin{lemma}[Grönwall]\label{lem:b3:ode:gronwall}
Misalkan $u \colon I \to \intco0\infty$ \omterm{def:b3:topology:continuity}{kontinu}, $t_0 \in
I$, dan andaikan\index{lema Grönwall}
\[
u(t) \leq \alpha + \beta\,\Bigl|\int_{t_0}^{t}u(s)\,\dd
s\Bigr| \qquad (t \in I)
\]
dengan $\alpha \geq 0$ dan $\beta > 0$. Maka $u(t) \leq
\alpha\,\eu^{\beta\abs{t - t_0}}$ pada $I$.
\end{lemma}

\begin{proof}
Untuk $t \geq t_0$: misalkan $v(t) = \alpha +
\beta\int_{t_0}^tu(s)\dd s$, sehingga $u \leq v$, $v' = \beta u
\leq \beta v$, dan $(v\eu^{-\beta(t - t_0)})' \leq 0$: jadi $v(t)
\leq v(t_0)\eu^{\beta(t-t_0)} = \alpha\eu^{\beta(t - t_0)}$.
Sedangkan untuk $t \leq t_0$, terapkan hal yang sama pada $\tilde u(t) = u(2t_0 -
t)$.
\end{proof}

\begin{corollary}[Ketunggalan dan kebergantungan kontinu]
\label{cor:b3:ode:uniqueness}
Di bawah hipotesis \cref{thm:b3:ode:cauchylipschitz},
dua solusi $x' = f(t,x)$ yang bersesuaian di satu titik bersesuaian
pada selang pendefinisian bersamanya. Secara kuantitatif, jika
$x, y$ dua solusi yang grafiknya berada di daerah yang $f$-nya bersifat
Lipschitz-$L$ terhadap $x$, maka
\[
\norm{x(t) - y(t)} \leq \norm{x(t_0) -
y(t_0)}\,\eu^{L\abs{t - t_0}} .
\]
\end{corollary}

\begin{proof}
Untuk taksirannya: $u = \norm{x - y}$ memenuhi $u(t) \leq u(t_0)
+ L\abs{\int_{t_0}^tu}$ (kurangkanlah persamaan integralnya);
lalu Grönwall. Untuk ketunggalan globalnya: himpunan kesesuaiannya $\{t : x(t)
= y(t)\}$ bersifat tertutup di selang bersamanya, tak kosong, dan
terbuka --- sebab di sekitar sembarang titik kesesuaiannya, selimutilah sepotong \omterm{def:b3:topology:compact}{kompak}
grafik bersamanya oleh berhingga banyak kotak Lipschitz lalu terapkan
taksirannya dengan $u(t_1) = 0$ pada masing-masingnya: sehingga secara lokal $x \equiv y$.
Dan himpunan bagian sebuah selang yang tak kosong, terbuka, dan tertutup adalah selangnya sendiri.
\end{proof}

\section{Solusi maksimal}

\begin{theorem}[Solusi maksimal; pelolosan dari kompak]
\label{thm:b3:ode:maximal}
Andaikan $f$ \omterm{def:b3:topology:continuity}{kontinu} dan \omterm{def:b3:ode:lipschitz}{Lipschitz lokal} terhadap $x$.
\begin{enumerate}
  \item Setiap soal Cauchy mempunyai satu-satunya solusi
        \emph{maksimal} $x \colon \intoo{T_-}{T_+} \to \R^d$: sehingga setiap
        solusi lain yang melalui $(t_0, x_0)$ merupakan
        pembatasannya. Dan selangnya terbuka.
  \item (\emph{Pelolosan}) Untuk setiap \omterm{def:b3:topology:compact}{kompak} $K \subseteq U$
        ada $\varepsilon > 0$ sedemikian sehingga $(t, x(t)) \notin
        K$ untuk setiap $t \in \intoo{T_+ - \varepsilon}{T_+}$ (dan
        setangkup dengan itu di $T_-$): yakni \emph{grafik sebuah \omterm{thm:b3:ode:maximal}{solusi
        maksimal} akhirnya meninggalkan setiap himpunan bagian \omterm{def:b3:topology:compact}{kompak}
        $U$}. Khususnya,
        untuk $U = \R\times\R^d$ dan $T_+ < +\infty$:
        $\norm{x(t)} \to +\infty$ saat $t \to T_+^-$
        (yakni \emph{peledakan}).\index{solusi maksimal}\index{pelolosan
        dari kompak}\index{peledakan}
\end{enumerate}
\end{theorem}

\begin{proof}
(1) Misalkan $\mathcal S$ himpunan semua solusi yang melalui
$(t_0, x_0)$; maka menurut \cref{cor:b3:ode:uniqueness} dua di antaranya bersesuaian pada
irisan selangnya, sehingga keduanya merekat: pada $J =
\bigcup_{y \in \mathcal S}I_y$, definisikan $x(t) = y(t)$ untuk sembarang
$y$ yang terdefinisi di $t$: yakni sebuah solusi yang terdefinisi baik, jelas
maksimal dan tunggal. Sedangkan $J$ terbuka: sebab sebuah solusi yang terdefinisi di
ujungnya dapat diperpanjang lewat
\cref{thm:b3:ode:cauchylipschitz} di ujung itu.

(2) Andaikan klaimnya gagal di $T_+$: maka ada $t_n \to
T_+$ dengan $(t_n, x(t_n)) \in K$; dan perhatikan bahwa ini memaksa $T_+ <
\infty$ atau, jika $T_+ = \infty$, maka tak ada yang perlu dibuktikan (sebab $K$
terbatas terhadap waktunya). Jadi misalkan $T_+ < \infty$. Lewat kekompakannya: konstanta seragam $M, L, a,
b$ berhasil bagi setiap data Cauchy pada sebuah \omterm{def:b3:topology:topology}{persekitaran} $K$
--- secara konkret, selimutilah $K$ oleh berhingga banyak kotak $Q_i$ seperti pada
bukti teorema lokalnya lalu misalkan $T^* > 0$ minimum
waktu keberadaan yang bersesuaian: sehingga setiap data Cauchy di $K$
meluncurkan sebuah solusi yang hidup setidaknya $T^*$ melampaui waktu awalnya.
Lalu menerapkan ini di $(t_n, x(t_n))$ dengan $t_n > T_+ -
T^*/2$ memperluas $x$ melewati $T_+$ (sebab perluasannya bersesuaian dengan $x$
menurut ketunggalannya, lalu memperpanjangnya): yang bertentangan dengan
kemaksimalannya. Jadi grafiknya meninggalkan $K$ secara pasti sebelum
$T_+$. Sedangkan untuk $U = \R\times\R^d$: jika $\norm{x(t)}\not\to\infty$,
maka sebuah barisan $t_n \to T_+$ menjaga $(t_n, x(t_n))$ di \omterm{def:b3:topology:compact}{kompak}
$[t_0, T_+]\times\bar B(0, R)$: yang terkecualikan.
\end{proof}

\begin{corollary}[Keberadaan global di bawah pertumbuhan linear]
\label{cor:b3:ode:globallinear}
Jika $U = I\times\R^d$ (dengan $I$ selang terbuka) dan $\norm{f(t, x)}
\leq \alpha(t)\norm x + \beta(t)$ dengan $\alpha, \beta$ yang
\omterm{def:b3:topology:continuity}{kontinu}, maka setiap \omterm{thm:b3:ode:maximal}{solusi maksimalnya} terdefinisi pada seluruh
$I$.
\end{corollary}

\begin{proof}
Pada sebuah \omterm{def:b3:topology:compact}{kompak} $[t_0, T] \subseteq I$: $\norm{x(t)} \leq
\norm{x_0} + \int_{t_0}^t(\alpha\norm x + \beta)$, sehingga menurut
Grönwall (dengan $\alpha, \beta$ terbatas oleh $A, B$ di sana)
$\norm{x(t)} \leq (\norm{x_0} + B(T - t_0))\eu^{A(T - t_0)}$:
jadi terbatas. Lalu jika $T_+ < \sup I$, grafiknya tinggal di sebuah \omterm{def:b3:topology:compact}{kompak}
$I\times\R^d$ di dekat $T_+$: yang bertentangan dengan pelolosannya
(\cref{thm:b3:ode:maximal}).
\end{proof}

\section{Sistem linear}

Sepanjang bagian ini $A \colon I \to M_d(\R)$ dan $b
\colon I \to \R^d$ bersifat \omterm{def:b3:topology:continuity}{kontinu}; sedangkan sistemnya adalah $x' = A(t)x
+ b(t)$ --- yakni pertumbuhan linear: sehingga semua \omterm{thm:b3:ode:maximal}{solusi maksimalnya} hidup pada seluruh
$I$ (\cref{cor:b3:ode:globallinear}).

\begin{theorem}[Struktur]\label{thm:b3:ode:linearstructure}
Solusi sistem homogen $x' = A(t)x$ membentuk ruang
vektor berdimensi $d$, yakni $S_H$; dan untuk setiap $t_0$,
penilaian $x \mapsto x(t_0)$ merupakan isomorfisma $S_H \to
\R^d$. Sedangkan \emph{resolven}\index{resolven sebuah sistem
linear} $R(t, s) \in GL_d(\R)$, yang didefinisikan lewat: $t \mapsto R(t,
s)v$ adalah solusi yang bernilai $v$ di $s$, memenuhi
\[
R(s,s) = I,\quad R(t, u)R(u, s) = R(t, s),\quad
\partial_tR(t,s) = A(t)R(t,s),
\]
dan soal tak homogennya diselesaikan \emph{rumus
Duhamel}\index{rumus Duhamel}:
\[
x(t) = R(t, t_0)\,x_0 + \int_{t_0}^{t}R(t, s)\,b(s)\,\dd s .
\]
Akhirnya \emph{Wronskian}\index{Wronskian} $w(t) = \det
R(t, s)$ menuruti rumus Liouville $w' =
\operatorname{tr}A(t)\,w$, sehingga $w(t) =
\exp\bigl(\int_s^t\operatorname{tr}A\bigr) > 0$.
\end{theorem}

\begin{proof}
Kelinearan persamaannya membuat solusinya menjadi ruang vektor;
sedangkan penilaiannya linear, injektif (menurut ketunggalannya: sebab sebuah solusi
yang lenyap di $t_0$ bernilai $\equiv 0$) dan surjektif (menurut keberadaannya):
jadi berdimensi $d$. Sifat \omterm{thm:b3:ode:linearstructure}{resolvennya} menyatakan ulang ketunggalannya
(sebab kedua ruas setiap identitasnya menyelesaikan soal Cauchy yang sama);
sedangkan keterbalikannya dari $R(s,t)R(t,s) = I$. Untuk Duhamel: turunkanlah
rumusnya --- $x'(t) = A(t)R(t,t_0)x_0 + R(t,t)b(t) +
\int_{t_0}^tA(t)R(t,s)b(s)\dd s = A(t)x(t) + b(t)$
(dan pendiferensialan di bawah integralnya sah: sebab
integrannya $\mathcal C^1$ terhadap $t$ dengan turunan yang \omterm{def:b3:topology:continuity}{kontinu}
terhadap $(t,s)$; atau periksalah lewat persamaan integralnya). Untuk Liouville:
$w(t + h) = \det\bigl(R(t+h, t)\bigr)w(t)$ dan $R(t + h, t) =
I + hA(t) + o(h)$ (dari persamaan integralnya), sehingga $\det = 1
+ h\operatorname{tr}A(t) + o(h)$ (lewat uraian $\det$ di
$I$): jadi $w'(t) = \operatorname{tr}A(t)\,w(t)$; lalu integralkan
persamaan diferensial linear skalarnya.
\end{proof}

\begin{theorem}[Eksponensial matriks]\label{thm:b3:ode:matrixexp}
Untuk $A \in M_d(\C)$, deret $\eu^{A} =
\sum_{n\geq0}\frac{A^n}{n!}$ konvergen (secara mutlak, dalam sembarang
norma submultiplikatif), $\eu^{A+B} = \eu^A\eu^B$ setiap kali
$AB = BA$, dan $t \mapsto \eu^{tA}$ merupakan \omterm{thm:b3:ode:linearstructure}{resolven}
sistem berkoefisien konstannya: $R(t,s) = \eu^{(t-s)A}$; dan ia bersifat $\mathcal
C^\infty$ dengan $\frac{\dd}{\dd t}\eu^{tA} = A\eu^{tA}$.
Lebih lanjut: jika $\operatorname{Re}\lambda < -\alpha < 0$ untuk
setiap nilai eigen $\lambda$ milik $A$, maka $\vertiii{\eu^{tA}}
\leq C\eu^{-\alpha t}$ untuk $t \geq 0$.\index{eksponensial
matriks}
\end{theorem}

\begin{proof}
Untuk kekonvergenannya: $\vertiii{A^n/n!} \leq \vertiii A^n/n!$, yang terjumlahkan
(\cref{exo:b3:complete:1}(b) pada aljabar Banach
$M_d$). Untuk $A, B$ yang komutatif: hasil kali Cauchy kedua
deret yang konvergen mutlak itu tersusun ulang, lewat teorema
binomialnya (yang sahih bila $AB = BA$), menjadi
$\sum_n\frac{(A+B)^n}{n!}$. Untuk keterturunannya, secara langsung:
$\frac{\eu^{(t+h)A} - \eu^{tA}}h = \eu^{tA}\frac{\eu^{hA} -
I}{h} \to \eu^{tA}A$ sebab $\norm{\frac{\eu^{hA} - I}h - A}
\leq \sum_{n\geq2}\frac{\abs h^{n-1}\vertiii A^n}{n!} =
O(h)$. Karena itu $t \mapsto \eu^{(t - s)A}v$ menyelesaikan soal
Cauchy yang mendefinisikan $R(t,s)v$. Untuk batas spektralnya: lewat \omterm{thm:b3:modules:jordan}{bentuk Jordan}
(\cref{thm:b3:modules:jordan}), $A = P(D + N)P^{-1}$ dengan $D$
diagonal yang membawa nilai eigennya, $N$ nilpoten, dan $DN =
ND$. Maka $\eu^{tA} = P\,\eu^{tD}\eu^{tN}P^{-1}$ dengan
$\vertiii{\eu^{tD}} \leq \eu^{-(\alpha + \delta)t}$ untuk suatu
$\delta > 0$ (dan $t \geq 0$) sedangkan $\eu^{tN}$ polinomial terhadap $t$
(menurut kenilpotenannya): sehingga hasil kalinya $\leq C\eu^{-\alpha t}$
(sebab polinomialnya dikalahkan $\eu^{-\delta t}$).
\end{proof}

\begin{example}[Bidangnya, yang digolongkan]
\label{ex:b3:ode:planeclassification}
Untuk $x' = Ax$ dengan $A \in M_2(\R)$ yang terbalikkan, potret
fasenya di dekat $0$ ditentukan oleh $\tau = \operatorname{tr}A$
dan $\delta = \det A$, lewat nilai eigennya $\lambda_\pm
= \frac{\tau \pm \sqrt{\tau^2 - 4\delta}}2$:
\begin{itemize}
  \item $\delta < 0$: nilai eigen real bertanda berlawanan ---
        yakni sebuah \emph{pelana}; dua \omterm{def:b3:topology:pathconnected}{lintasan} masuk, dua keluar,
        dan sisanya melintas. Selalu tak \omterm{def:b3:ode:stability}{stabil}.
  \item $\delta > 0$ dan $\tau^2 \geq 4\delta$: nilai eigen real
        bertanda sama (yakni $= \operatorname{sign}\tau$) --- yakni sebuah
        \emph{simpul}, yang \omterm{def:b3:ode:stability}{stabil} jika dan hanya jika $\tau < 0$; dan \omterm{def:b3:topology:pathconnected}{lintasannya}
        menyinggung arah eigen yang lambat.
  \item $\delta > 0$, $\tau^2 < 4\delta$, dan $\tau \neq 0$:
        nilai eigen sekawan kompleks $\frac\tau2 \pm
        \iu\omega$ --- yakni sebuah \emph{spiral} (atau fokus), yang \omterm{def:b3:ode:stability}{stabil} jika dan hanya jika
        $\tau < 0$; sedangkan solusinya berupa
        $\eu^{\tau t/2}\times$ rotasi berperiode
        $\frac{2\pi}\omega$.
  \item $\tau = 0$ dan $\delta > 0$: nilai eigen murni
        imajiner --- yakni sebuah \emph{pusat}: dengan \omterm{def:b3:groups:action}{orbit} tertutup
        (elips), yakni kestabilan tanpa kestabilan asimtotik,
        tepat perbatasan yang tak dapat diputuskan
        \cref{thm:b3:ode:linearization} bagi sistem
        taklinear (sebab kesetimbangan bawah bandulnya,
        \cref{pb:b3:ode:1}, duduk di sini).
\end{itemize}
Sedangkan parabola batasnya $\tau^2 = 4\delta$ membawa simpul yang
merosot (yakni blok Jordan: dengan \omterm{def:b3:topology:pathconnected}{lintasan} berarah singgung
tunggal). Jadi segalanya terbaca dari dua bilangan ---
itulah mengapa refleks pertama sebelum menggambar sebuah potret fase
bidang adalah menghitung $\operatorname{tr}$ dan $\det$; misalnya
$A = \bigl(\begin{smallmatrix}0 & 1\\ -1 &
-c\end{smallmatrix}\bigr)$ (yakni osilator teredam): $\delta = 1 >
0$ dan $\tau = -c$: sehingga spiral \omterm{def:b3:ode:stability}{stabil} untuk $0 < c < 2$, dan simpul \omterm{def:b3:ode:stability}{stabil}
untuk $c \geq 2$ --- yakni redaman kurang lawan redaman lebih, dalam satu
pandangan.
\end{example}

\section{Aliran, kesetimbangan, kestabilan}

Kini tinjaulah persamaan \emph{otonom} $x' = F(x)$, dengan $F
\colon \Omega \to \R^d$ yang \omterm{def:b3:ode:lipschitz}{Lipschitz lokal} pada $\Omega
\subseteq \R^d$ yang terbuka. Tulislah $\varphi_t(x_0) = x(t)$ untuk
\omterm{thm:b3:ode:maximal}{solusi maksimal} dengan $x(0) = x_0$ (yakni
\emph{aliran}\index{aliran sebuah medan vektor} itu); sedangkan keotonomannya memberikan
sifat grupnya $\varphi_{t+s} =
\varphi_t\circ\varphi_s$ di tempat ia terdefinisi (sebab kedua ruasnya menyelesaikan
soal yang sama pada waktu $s$).

\begin{definition}\label{def:b3:ode:stability}
Sebuah \emph{kesetimbangan} adalah titik $\bar x$ dengan $F(\bar x) =
0$ (sehingga $\varphi_t(\bar x) = \bar x$). Ia disebut
\emph{stabil}\index{kestabilan (Lyapunov)} jika untuk setiap
$\varepsilon > 0$ ada $\delta > 0$ sedemikian sehingga
$\norm{x_0 - \bar x} < \delta$ mengakibatkan solusinya
ada untuk setiap $t \geq 0$ dengan $\norm{\varphi_t(x_0) - \bar
x} < \varepsilon$; dan \emph{stabil asimtotik} jika lebih lanjut
$\varphi_t(x_0) \to \bar x$ untuk setiap $x_0$ di dekat $\bar x$.
\end{definition}

\begin{theorem}[Fungsi Lyapunov]\label{thm:b3:ode:lyapunov}
Misalkan $\bar x$ sebuah kesetimbangan dan $V \colon \mathcal V \to
\R$ sebuah fungsi $\mathcal C^1$ pada sebuah \omterm{def:b3:topology:topology}{persekitaran} $\bar x$
dengan:\index{fungsi Lyapunov}
\[
V(\bar x) = 0,\qquad V(x) > 0 \text{ untuk } x \neq \bar x,
\qquad \dot V(x) = \nabla V(x)\cdot F(x) \leq 0 .
\]
Maka $\bar x$ bersifat \omterm{def:b3:ode:stability}{stabil}. Dan jika lebih lanjut $\dot V < 0$ di luar $\bar
x$, maka $\bar x$ \omterm{def:b3:ode:stability}{stabil} asimtotik.
\end{theorem}

\begin{proof}
Sepanjang sebuah solusi, $\frac{\dd}{\dd t}V(x(t)) = \dot V(x(t))
\leq 0$: sehingga $V$ menurun. Diberikan $\varepsilon$ (yang cukup kecil
sehingga $\bar B(\bar x, \varepsilon) \subseteq \mathcal V$), misalkan
$m = \min\{V(x) : \norm{x - \bar x} = \varepsilon\} > 0$
(lewat kekompakan dan kepositifannya) lalu pilihlah $\delta < \varepsilon$
dengan $V < m$ pada $B(\bar x, \delta)$ (lewat \omterm{def:b3:topology:continuity}{kekontinuannya}). Maka sebuah solusi
yang bermula di $B(\bar x, \delta)$ mempunyai $V(x(t)) < m$ untuk setiap
waktu berikutnya, sehingga ia tak pernah dapat mencapai bola $\norm{x -
\bar x} = \varepsilon$ (yang $V \geq m$ padanya): jadi ia tinggal di
bolanya --- dan lalu ada untuk setiap $t \geq 0$: sebab solusinya
tetap di $\bar B$ yang \omterm{def:b3:topology:compact}{kompak}, sehingga
\cref{thm:b3:ode:maximal}(2) (yakni \omterm{thm:b3:ode:maximal}{pelolosan dari kompak}) memaksa
$T_+ = +\infty$. Jadi kestabilannya.

Untuk kasus asimtotiknya: misalkan $x(t)$ bermula di $B(\bar x, \delta)$;
maka $V(x(t))$ menurun ke suatu $c \geq 0$. Seandainya $c > 0$: maka
\omterm{def:b3:topology:pathconnected}{lintasannya} tinggal di $K = \{x \in \bar B(\bar x,
\varepsilon): V(x) \geq c\}$, yakni himpunan \omterm{def:b3:topology:compact}{kompak} yang mengecualikan sebuah
\omterm{def:b3:topology:topology}{persekitaran} $\bar x$ (sebab $V$ \omterm{def:b3:topology:continuity}{kontinu} dengan $V(\bar x)
= 0 < c$). Lalu pada $K$, fungsi $\dot V$ bersifat \omterm{def:b3:topology:continuity}{kontinu} dan
negatif sejati, sehingga $\mu = \max_K\dot V < 0$
(lewat kekompakannya); maka $V(x(t)) \leq V(x(0)) + \mu t \to
-\infty$: yang mustahil, sebab $V \geq 0$. Jadi
$c = 0$, dan $x(t) \to \bar x$ (sebab titik yang berjarak $\geq
\rho$ dari $\bar x$ di dalam bolanya mempunyai $V \geq m_\rho >
0$).
\end{proof}

\begin{theorem}[Kestabilan lewat pelinearan]
\label{thm:b3:ode:linearization}
Misalkan $F$ bersifat $\mathcal C^1$, $F(\bar x) = 0$, dan $A = DF(\bar x)$.
Jika setiap nilai eigen $A$ mempunyai $\operatorname{Re}\lambda <
0$, maka $\bar x$ bersifat \omterm{def:b3:ode:stability}{stabil} asimtotik.\index{kestabilan lewat
pelinearan}
\end{theorem}

\begin{proof}
Geserlah $\bar x$ ke $0$ lalu tulislah $F(x) = Ax + g(x)$ dengan
$g(x) = o(\norm x)$ (menurut keterturunan $\mathcal C^1$-nya). Pilihlah
$\alpha > 0$ dengan $\vertiii{\eu^{tA}} \leq C\eu^{-\alpha t}$
(untuk $t \geq 0$; \cref{thm:b3:ode:matrixexp}) dan $r > 0$ dengan
$\norm{g(x)} \leq \frac{\alpha}{2C}\norm x$ untuk $\norm x \leq
r$. Lalu Duhamel dengan $b(s) = g(x(s))$:
\[
x(t) = \eu^{tA}x_0 + \int_0^t\eu^{(t-s)A}g(x(s))\,\dd s,
\]
yang sahih selama $\norm{x(s)} \leq r$. Maka $u(t) =
\eu^{\alpha t}\norm{x(t)}$ memenuhi
\[
u(t) \leq C\norm{x_0} +
\int_0^{t}C\,\frac{\alpha}{2C}\,u(s)\,\dd s ,
\]
sehingga Grönwall memberikan $u(t) \leq C\norm{x_0}\eu^{\alpha t/2}$,
yakni $\norm{x(t)} \leq C\norm{x_0}\eu^{-\alpha t/2}$. Lalu jika
$\norm{x_0} < r/C$, batas a priorinya menjaga $\norm{x(t)} <
r$ untuk setiap $t$ (lewat argumen \omterm{def:b3:topology:continuity}{kekontinuan} atau pengangkatan: sebab himpunan
waktu yang $\norm x \leq r$ padanya bersifat terbuka sekaligus tertutup di
$\intco0{T_+}$ mengingat taksiran sejatinya), sehingga solusinya
global, dan ia konvergen ke $0$ secara eksponensial: yakni kestabilan
asimtotiknya.
\end{proof}

\begin{method}\label{met:b3:ode:strategy}
Menghadapi sebuah persamaan diferensial: (1) untuk keberadaan dan ketunggalannya --- periksalah
kelipschitzan lokalnya (biasanya $\mathcal C^1$); (2) untuk keglobalannya --- lewat pertumbuhan
linear, keterbatasan, atau sebuah \omterm{def:b3:topology:compact}{kompak} invarian lewat \omterm{thm:b3:ode:lyapunov}{fungsi Lyapunov}
atau integral pertama; dan jika gagal, curigailah peledakan
lalu ujilah pada karikatur skalarnya $x' = x^2$; (3) untuk sistem
linearnya --- lewat \omterm{thm:b3:ode:linearstructure}{resolven}, Duhamel, dan untuk koefisien
konstannya lewat struktur eigen $A$; (4) untuk pertanyaan
kualitatifnya --- lewat kesetimbangan, pelinearan, lalu perburuan \omterm{thm:b3:ode:lyapunov}{fungsi
Lyapunov} (yakni energinya, bila sistemnya mekanis) atau sebuah integral
pertama yang himpunan arasnya menjebak \omterm{def:b3:topology:pathconnected}{lintasannya}. Sedangkan soal akhir
pekannya menjalankan seluruh metode itu pada bandulnya.
\end{method}

\begin{omfigure}
\begin{tikzpicture}
\begin{axis}[omaxis, width=11cm, height=6cm,
    xmin=-7.2, xmax=7.2, ymin=-2.9, ymax=2.9,
    xtick={-6.2832, -3.1416, 0, 3.1416, 6.2832},
    xticklabels={$-2\pi$, $-\pi$, $0$, $\pi$, $2\pi$},
    ytick={-2,0,2}]
  % oscillations: E<1 : y = +-sqrt(2(E+cos x)), E = -0.3, 0.5
  \foreach \E in {-0.3, 0.5}{
    \addplot[omDef, samples=200, domain=-7.2:7.2]
      {sqrt(max(0, 2*(\E + cos(deg(x)))))};
    \addplot[omDef, samples=200, domain=-7.2:7.2]
      {-sqrt(max(0, 2*(\E + cos(deg(x)))))};
  }
  % separatrix E = 1
  \addplot[omThm, thick, samples=300, domain=-7.2:7.2]
    {sqrt(max(0, 2*(1 + cos(deg(x)))))};
  \addplot[omThm, thick, samples=300, domain=-7.2:7.2]
    {-sqrt(max(0, 2*(1 + cos(deg(x)))))};
  % rotations E > 1
  \foreach \E in {1.6, 2.4}{
    \addplot[omProp, samples=200, domain=-7.2:7.2]
      {sqrt(2*(\E + cos(deg(x))))};
    \addplot[omProp, samples=200, domain=-7.2:7.2]
      {-sqrt(2*(\E + cos(deg(x))))};
  }
\end{axis}
\end{tikzpicture}

{\small Potret fase bandul $x'' = -\sin x$ pada
bidang $(x, x')$: yakni kurva aras energinya $E = \frac{x'^2}2
- \cos x$. Kurva tertutupnya (biru): osilasi, dengan $E < 1$;
kurva berjalannya (jingga): rotasi penuh, dengan $E > 1$; dan di antara
keduanya \emph{separatriksnya} (merah), $E = 1$, yang menghubungkan
kesetimbangan tak \omterm{def:b3:ode:stability}{stabilnya} $(\pm\pi, 0)$. Sedangkan soal akhir pekannya membuktikan
segala yang disarankan gambar ini.}
\end{omfigure}

\section{Latihan}

\begin{exercise}[$\star$]\label{exo:b3:ode:1}
Selesaikan secara eksplisit lalu tentukan selang maksimalnya:
(a) $x' = x^2$, $x(0) = 1$; (b) $x' = 1 + x^2$, $x(0) = 0$;
(c) $x' = x(1-x)$, $x(0) = \frac12$. Lalu damaikan setiap jawabannya
dengan \cref{thm:b3:ode:maximal}(2) dan
\cref{cor:b3:ode:globallinear}.
\end{exercise}

\begin{exercise}[$\star$]\label{exo:b3:ode:2}
Misalkan $x, y$ menyelesaikan $x' = f(t,x)$ dengan $f$ yang Lipschitz-$L$
secara global terhadap $x$ pada $\R\times\R^d$.
(a) Buktikan $\norm{x(t) - y(t)} \leq \norm{x(0) -
y(0)}\eu^{L\abs t}$, lalu tunjukkan lewat contoh (yang linear!) bahwa
faktor $\eu^{L\abs t}$ tercapai.
(b) Turunkan bahwa pemetaan \omterm{ex:b3:ode:planeclassification}{aliran} $x_0 \mapsto x(t; x_0)$ bersifat
\omterm{def:b3:topology:continuity}{kontinu}, bahkan Lipschitz pada himpunan terbatas.
\end{exercise}

\begin{exercise}[$\star\star$]\label{exo:b3:ode:3}
Tunjukkan bahwa masing-masing berikut ini bersolusi maksimal yang semuanya
global pada $\R$, sambil mengutip teorema yang tepat:
(a) $x' = \sin(tx)$;
(b) $x' = \frac{t\,x}{1 + x^2}$;
(c) $x'' + q(t)x = 0$ dengan $q$ yang \omterm{def:b3:topology:continuity}{kontinu} (ubahlah menjadi
sistem orde pertama);
(d) $x' = A(t)x$ dengan $A$ yang \omterm{def:b3:topology:continuity}{kontinu} dan terbatas --- lalu berikan
batas Grönwall bagi $\norm{x(t)}$.
\end{exercise}

\begin{exercise}[$\star\star$]\label{exo:b3:ode:4}
(a) Hitunglah $\eu^{tA}$ untuk
$A = \bigl(\begin{smallmatrix}0 & -1\\ 1 &
0\end{smallmatrix}\bigr)$,
$\bigl(\begin{smallmatrix}\lambda & 1\\ 0 &
\lambda\end{smallmatrix}\bigr)$, dan
$\bigl(\begin{smallmatrix}0 & 1\\ 1 &
0\end{smallmatrix}\bigr)$.
(b) Selesaikan osilator yang digerakkan $x'' + x = \cos(\omega t)$
lewat Duhamel (dalam bentuk sistemnya), untuk $\omega \neq 1$ dan $\omega =
1$: sehingga resonansinya muncul sebagai suku sekular $t\sin t$.
\end{exercise}

\begin{exercise}[$\star\star$]\label{exo:b3:ode:5}
Untuk persamaan skalar $x'' + p(t)x' + q(t)x = 0$:
(a) Tunjukkan bahwa \omterm{thm:b3:ode:linearstructure}{Wronskian} $w = x_1x_2' - x_1'x_2$ dua
solusinya memenuhi $w' = -p\,w$ (menurut Abel), lalu turunkan bahwa dua
solusi yang $w \neq 0$-nya di suatu tempat membentuk sebuah basis.
(b) Diberikan satu solusi $x_1$ yang tak lenyap, carilah solusi
umumnya lewat \emph{penurunan orde}: tetapkan $x_2 = x_1\int
\frac{\eu^{-\int p}}{x_1^2}$ lalu periksalah. Lalu terapkan pada $t^2x'' -
2x = 0$ di $\intoo0{+\infty}$ dengan $x_1(t) = t^2$.
\end{exercise}

\begin{exercise}[$\star\star$]\label{exo:b3:ode:6}
(Logistik) Untuk $x' = x(1 - x)$: tentukan semua kesetimbangannya dan
kestabilannya (lewat \cref{thm:b3:ode:linearization} dan
secara langsung); lalu tunjukkan bahwa setiap solusi dengan $x(0) \in \intoo01$ bersifat
naik dan global, dengan limit $0$ dan $1$ di $\mp\infty$;
lalu selesaikan secara eksplisit untuk membenarkannya. Lalu tunjukkan secara lebih umum bahwa
solusi otonom skalarnya bersifat monoton, lalu simpulkan:
tak ada solusi periodik tak konstan dalam dimensi $1$.
\end{exercise}

\begin{exercise}[$\star\star$]\label{exo:b3:ode:7}
(Integral pertama) Misalkan $H \colon \Omega \to \R$ bersifat $\mathcal
C^1$ lalu tinjaulah sistem Hamilton bidang $x' =
\partial_yH$, $y' = -\partial_xH$.
(a) Tunjukkan bahwa $H$ konstan sepanjang solusinya.
(b) Untuk $H = \frac{y^2}2 + \frac{x^4}4$: tunjukkan bahwa semua solusinya
global dan terbatas, dan bahwa titik asalnya \omterm{def:b3:ode:stability}{stabil}
(lewat Lyapunov: yakni $H$) walaupun pelinearannya
($\bigl(\begin{smallmatrix}0&1\\0&0\end{smallmatrix}\bigr)$)
\emph{tak} \omterm{def:b3:ode:stability}{stabil} asimtotik: sehingga pelinearannya dapat tak menyimpulkan.
\end{exercise}

\begin{exercise}[$\star\star\star$]\label{exo:b3:ode:8}
(Bandul teredam) $x'' + cx' + \sin x = 0$ dengan $c > 0$; dan sistemnya:
$x' = y$, $y' = -\sin x - cy$.
(a) Tunjukkan bahwa $V(x, y) = \frac{y^2}2 + 1 - \cos x$ memenuhi $\dot
V = -cy^2 \leq 0$: sehingga titik asalnya \omterm{def:b3:ode:stability}{stabil}.
(b) Sedangkan $\dot V$ lenyap pada seluruh sumbu $y = 0$: sehingga kriteria
sejati Lyapunov gagal. Buktikanlah kestabilan asimtotiknya tetap,
lewat pelinearan (\cref{thm:b3:ode:linearization}): hitunglah
nilai eigen matriks terlinearkannya di $(0,0)$ lalu periksalah
$\operatorname{Re} < 0$ untuk setiap $c > 0$.
(c) Apa yang terjadi pada kesetimbangan $(\pi, 0)$? Hitunglah
pelinearannya lalu simpulkan (sebab satu nilai eigennya positif:
yakni ketakstabilan --- Anda boleh memakai pernyataan ketakstabilannya
secara tak formal atau menghasilkan sebuah solusi yang lolos secara eksplisit dari
sistem linearnya).
\end{exercise}

\begin{exercise}[$\star\star$]\label{exo:b3:ode:9}
(Perbatasan ketunggalannya) Untuk $\alpha \in \intoo01$, tunjukkan bahwa
soal $x' = \abs x^{\alpha}$ dengan $x(0) = 0$ mempunyai tak berhingga
banyak solusi (sesuaikan \cref{pb:b3:complete:1}, Bagian III).
Lalu tunjukkan sebaliknya bahwa untuk $\alpha = 1$ (yakni $x' = \abs
x$) solusi yang melalui $0$ bersifat tunggal, lalu kenali
dengan tepat hipotesis mana pada
\cref{thm:b3:ode:cauchylipschitz} yang membedakan kedua
kasusnya.
\end{exercise}

\begin{exercise}[$\star\star\star$]\label{exo:b3:ode:10}
(Batas a priori menjebak solusinya) Misalkan $F \colon \R^d \to \R^d$
\omterm{def:b3:ode:lipschitz}{Lipschitz lokal} dengan $\langle F(x), x\rangle \leq 0$
setiap kali $\norm x \geq R$.
(a) Tunjukkan bahwa bola tertutup $\bar B(0, R)$ bersifat invarian
secara positif: yakni solusi yang bermula di dalamnya tinggal di dalamnya untuk $t \geq
0$. \emph{(Jika $\norm{x(t_2)} > R$, tinjaulah waktu terakhir
$t_1 < t_2$ dengan $\norm{x(t_1)} = R$ lalu pelajarilah
$\frac{\dd}{\dd t}\norm{x(t)}^2$ pada $\intcc{t_1}{t_2}$.)}
(b) Turunkan keberadaan maju yang global bagi data di bolanya.
Lalu tanganilah sistem gradiennya $x' = -\nabla G(x)$ dengan $G \in
\mathcal C^2$ dan $G(x) \to +\infty$ saat $\norm x \to
\infty$: tunjukkan bahwa $G$ menurun sepanjang solusinya, bahwa setiap
solusinya tinggal di himpunan subaras (yang terbatas) $\{G \leq
G(x_0)\}$, lalu simpulkan keberadaan maju yang global.
\end{exercise}

\begin{exercise}[$\star\star$]\label{exo:b3:ode:11}
(Peledakan lewat pembandingan) Tinjaulah $x' = x^2 + t^2$ dengan $x(0) =
1$.
(a) Tunjukkan bahwa \omterm{thm:b3:ode:maximal}{solusi maksimalnya} ada pada suatu $\intco0{T_+}$
dengan $T_+ < \infty$: bandingkanlah dengan $y' = y^2$, $y(0) = 1$
\emph{(buktikanlah lema pembandingan yang Anda perlukan: bahwa jika $x' \geq
F(x)$ dan $y' = F(y)$ dengan $x(0) \geq y(0)$, maka $x \geq
y$ di tempat keduanya hidup)}, lalu turunkan $T_+ \leq 1$.
(b) Batasilah $T_+$ dari bawah: pada $\intcc01$, $x' \leq x^2 +
1$; lalu bandingkan dengan supersolusinya $z' = z^2 + 1$ dengan $z(0) =
1$, yang diselesaikan $z(t) = \tan\bigl(t + \frac\pi4\bigr)$, lalu
simpulkan $T_+ \geq \frac\pi4$.
(c) Rakitlah $\frac\pi4 \leq T_+ \leq 1$ lalu rumuskan
moralnya: bahwa pertumbuhan superlinear ruas kanannya adalah yang
membunuh keberadaan globalnya (dengan \cref{exo:b3:ode:3} sebagai
tandingannya), dengan perbatasannya berupa kekonvergenan
$\int^{\infty}\frac{\dd s}{F(s)}$.
\end{exercise}

\begin{exercise}[$\star\star\star$]\label{exo:b3:ode:12}
(Teorema pembandingan Sturm) Misalkan $q_1 \leq q_2$ \omterm{def:b3:topology:continuity}{kontinu}
pada sebuah selang $I$, dan misalkan $u \neq 0$ menyelesaikan $u''
+ q_1u = 0$ dan $v \neq 0$ menyelesaikan $v'' + q_2v = 0$.
(a) Tegakkanlah identitas \omterm{thm:b3:ode:linearstructure}{Wronskiannya}: bahwa dengan $W = uv' - u'v$,
$W' = (q_1 - q_2)\,uv$.
(b) (Sturm) Tunjukkan bahwa di antara dua akar berurutan $a < b$
milik $u$, entah $v$ lenyap di suatu tempat pada $\intoo ab$, entah
$q_1 = q_2$ dan $v \propto u$ di sana \emph{(andaikan $u > 0$
pada $\intoo ab$ dan $v > 0$ pula; lalu integralkan (a) dari $a$ ke
$b$ lalu periksalah tanda suku batasnya $W(a),
W(b)$)}.
(c) Turunkan: bahwa solusi $u'' + qu = 0$ dengan $q \geq m > 0$
lenyap setidaknya sekali pada setiap selang berpanjang
$\pi/\sqrt m$ (bandingkanlah dengan $v'' + mv = 0$); sedangkan solusi dengan
$q \leq 0$ lenyap paling banyak sekali pada $\R$. Lalu ujilah keduanya pada
$q \equiv \pm1$.
\end{exercise}

\section{Soal: bandulnya, yang terselesaikan sepenuhnya}

\begin{problem}[{Soal akhir pekan --- osilasi, rotasi,
separatriks, dan periodenya}]\label{pb:b3:ode:1}
Persamaan bandul $x'' = -\sin x$ --- sebagai sistem: $x' =
y$, $y' = -\sin x$ pada $\R^2$ --- merupakan lalat buah
dinamika: mudah ditulis, mustahil diselesaikan lewat rumus
dasar, namun sepenuhnya terpahami lewat metode kualitatifnya.
Misalkan $E(x, y) = \frac{y^2}2 - \cos x$ (yakni
\emph{energinya}).

\medskip
\noindent\textbf{Bagian I --- Struktur globalnya.}
\begin{enumerate}
  \item Tunjukkan bahwa semua \omterm{thm:b3:ode:maximal}{solusi maksimalnya} bersifat global
        (yakni terdefinisi pada $\R$): pakailah $\dot E = 0$ dan
        \cref{thm:b3:ode:maximal}. Untuk kesetimbangannya: $(k\pi, 0)$;
        lalu golongkan pelinearannya (bertipe pusat bagi $k$ yang genap,
        dan pelana bagi $k$ yang gasal).
  \item Tunjukkan bahwa \omterm{def:b3:topology:pathconnected}{lintasannya} termuat di himpunan aras
        $\{E = E_0\}$, lalu sketsakan atau uraikan menurut
        nilai $E_0 \in \intco{-1}{+\infty}$: $E_0 = -1$
        (kesetimbangannya), $-1 < E_0 < 1$ (kurva tertutup di sekitar
        $(2k\pi, 0)$), $E_0 = 1$ (separatriks yang melalui
        $(\pm\pi, 0)$), dan $E_0 > 1$ (grafik atas $x$:
        yakni rotasinya).
  \item Buktikan bahwa kesetimbangan bawahnya $(0,0)$ bersifat \omterm{def:b3:ode:stability}{stabil}
        tetapi \emph{tak} \omterm{def:b3:ode:stability}{stabil} asimtotik. \emph{(Pakailah Lyapunov
        dengan $V = E + 1$; sedangkan ketakasimtotikannya: sebab kekekalan energinya
        menjebak \omterm{def:b3:groups:action}{orbitnya} pada kurva aras yang jauh dari
        titik asalnya.)}
\end{enumerate}

\medskip
\noindent\textbf{Bagian II --- Osilasi dan periodenya.}
Tetapkan $-1 < E_0 < 1$ lalu tulislah $E_0 = -\cos a$ dengan $a \in
\intoo0\pi$ (yakni amplitudonya).
\begin{enumerate}[resume]
  \item Tunjukkan bahwa solusi dengan $x(0) = a$ dan $y(0) = 0$
        berosilasi: yakni $x(t) \in \intcc{-a}a$, dan \omterm{def:b3:groups:action}{orbitnya} berupa
        kurva tertutup $y^2 = 2(\cos x - \cos a)$. Lalu benarkanlah
        bahwa solusinya periodik: sebab \omterm{def:b3:groups:action}{orbitnya} kurva \omterm{def:b3:topology:compact}{kompak}
        tanpa kesetimbangan, yang dilalui dengan laju terbatas
        dari bawah --- jadikanlah ini sebuah argumen (sebab solusinya kembali
        ke titik awalnya dalam waktu berhingga, lalu ketunggalannya
        memaksa keperiodikannya).
  \item Tegakkanlah rumus periodenya
        \[
        T(a) = 4\int_0^{a}\frac{\dd x}{\sqrt{2(\cos x -
        \cos a)}}
        \]
        \emph{(pada seperempat \omterm{def:b3:groups:action}{orbitnya}, $y = \frac{\dd x}{\dd t} >
        0$ dan pisahkanlah variabelnya; lalu benarkan
        kekonvergenan tak wajarnya di $x = a$)}.
  \item (Osilasi kecil) Substitusikan $\sin\frac x2 =
        \sin\frac a2\,\sin\varphi$ lalu tunjukkan bahwa
        \[
        T(a) = 4\int_0^{\pi/2}\frac{\dd\varphi}{\sqrt{1 -
        k^2\sin^2\varphi}},
        \qquad k = \sin\frac a2
        \]
        (yakni sebuah integral eliptik \omterm{def:b3:complete:complete}{lengkap}). Lalu turunkan lewat kekonvergenan
        terdominasi bahwa $T(a) \to 2\pi$ saat $a \to 0^+$: yakni
        limit harmoniknya, yang tak bergantung pada amplitudonya ---
        yaitu keisokronan hampiran Galileo, beserta koreksi persisnya $T(a) = 2\pi\bigl(1 + \frac{k^2}4 +
        O(k^4)\bigr)$ (uraikanlah integrannya lalu integralkan
        suku demi suku, dan benarkanlah lewat kekonvergenan normalnya).
  \item Tunjukkan bahwa $T(a) \to +\infty$ saat $a \to \pi^-$
        \emph{(batasilah integrannya dari bawah di dekat $\varphi =
        \frac\pi2$ saat $k \to 1$, atau terapkan kekonvergenan
        monotonnya)}: sehingga saat mendekati separatriksnya,
        bandulnya melambat tanpa batas.
\end{enumerate}

\medskip
\noindent\textbf{Bagian III --- Separatriksnya.}
\begin{enumerate}[resume]
  \item Untuk $E_0 = 1$, $y = 2\cos\frac x2$ pada cabang
        atasnya: pisahkanlah variabelnya lalu integralkan untuk mencari
        solusi eksplisitnya
        \[
        x(t) = 4\arctan\bigl(\eu^{t}\bigr) - \pi
        \]
        (dengan $x(0) = 0$ dan $y(0) = 2$). Lalu periksalah secara langsung bahwa
        ia menyelesaikan persamaan bandulnya, lalu hitunglah limitnya beserta
        limit $y(t)$ saat $t \to
        \pm\infty$.
  \item Simpulkan: bahwa \omterm{def:b3:groups:action}{orbit} separatriksnya menghubungkan pelana
        $(-\pi, 0)$ (saat $t \to -\infty$) ke pelana
        $(\pi, 0)$ (saat $t \to +\infty$) tetapi tak mencapai satu pun
        dalam waktu berhingga --- yang konsisten dengan ketunggalannya (mengapa
        mencapai sebuah pelana dalam waktu berhingga akan bertentangan dengan
        \cref{cor:b3:ode:uniqueness}?).
\end{enumerate}

\medskip
\noindent\textbf{Bagian IV --- Rotasi, dan gambaran
\omterm{def:b3:complete:complete}{lengkapnya}.}
\begin{enumerate}[resume]
  \item Untuk $E_0 > 1$: tunjukkan bahwa $y$ tak pernah lenyap, $x$ bersifat
        monoton sejati dan global dengan $x(t) \to
        \pm\infty$, dan $t \mapsto y(t)$ bersifat periodik dengan
        periode
        \[
        \tau(E_0) = \int_{-\pi}^{\pi}
        \frac{\dd x}{\sqrt{2(E_0 + \cos x)}} .
        \]
  \item Rakitlah potret fase yang \omterm{def:b3:complete:complete}{lengkap} (yakni gambar bab ini)
        beserta pembenaran penuh setiap gejalanya, lalu
        tulislah ringkasan sepuluh baris atas metodenya: yakni energi,
        himpunan aras, kekompakan, dan ketunggalan --- bagaimana setiap
        teorema bab ini masuk. Lalu di manakah kita pernah memerlukan sebuah rumus bagi solusi umumnya?
\end{enumerate}

\medskip
\noindent\textbf{Bagian V --- Fungsi periodenya di bawah
mikroskop.}
\begin{enumerate}[resume]
  \item Buktikan momen Wallisnya
        \[
        W_n = \int_0^{\pi/2}\sin^{2n}\varphi\,\dd\varphi
        = \frac\pi2\cdot\frac{(2n)!}{4^n\,(n!)^2}
        \]
        lewat induksi (dengan mengintegralkannya secara parsial), lalu uraikan
        integran pertanyaan 6 lewat deret binomialnya, lalu
        benarkanlah pengintegralan suku demi sukunya untuk memperoleh deret
        \omterm{def:b3:complete:complete}{lengkapnya}
        \[
        T(a) = 2\pi\sum_{n\geq0}
        \Bigl(\frac{(2n)!}{4^n(n!)^2}\Bigr)^{\!2}k^{2n}
        = 2\pi\Bigl(1 + \frac{k^2}4 + \frac{9k^4}{64}
        + O(k^6)\Bigr),
        \qquad k = \sin\frac a2 .
        \]
  \item Ubahlah ke amplitudonya:
        \[
        T(a) = 2\pi\Bigl(1 + \frac{a^2}{16} +
        \frac{11\,a^4}{3072} + O(a^6)\Bigr)
        \]
        \emph{(substitusikanlah uraian $\sin\frac a2$
        lalu kumpulkan)}. Sehingga keisokronannya gagal pada orde $a^2$, dan
        kegagalannya kini \omterm{def:b3:lebesgue:measurable}{terukur} sampai orde $a^4$.
  \item Tunjukkan bahwa $a \mapsto T(a)$ bersifat \omterm{def:b3:topology:continuity}{kontinu} dan
        naik sejati pada $\intoo0\pi$, lalu simpulkan
        bersama pertanyaan 6--7 bahwa $T$ merupakan bijeksi dari
        $\intoo0\pi$ ke $\intoo{2\pi}{+\infty}$: sehingga setiap
        periode superkritisnya diwujudkan tepat satu amplitudo.
  \item (Aritmetika pembuat jam) Sebuah bandul yang diatur pada
        amplitudo yang melenyap menjaga waktu yang \omterm{def:b3:rings:ideal}{ideal}; lalu tunjukkan bahwa jika dijalankan
        pada amplitudo $a$ ia tertinggal sebesar pecahan
        $\frac{a^2}{16} + O(a^4)$ dari waktu \omterm{def:b3:rings:ideal}{idealnya}, lalu hitunglah
        pergeserannya untuk $a = 0.2$ rad: yakni sekitar $216$ detik per
        hari. (Sedangkan pipi sikloid Huygens dan amplitudo tetap yang kecil
        pada rodagigi pelepasnya sama-sama merupakan jawaban atas bilangan ini.)
  \item Kembalilah ke periode rotasi $\tau$ pada pertanyaan
        10: tunjukkan bahwa $\tau$ turun sejati pada
        $\intoo1{+\infty}$, bahwa $\tau(E_0) \to +\infty$ saat
        $E_0 \to 1^+$ (lewat kekonvergenan monoton), dan bahwa
        $\sqrt{2E_0}\,\tau(E_0) \to 2\pi$ saat $E_0 \to
        +\infty$ (lewat kekonvergenan terdominasi): sehingga pusaran cepat
        secara asimtotik merupakan rotasi bebas berlaju sudut
        $\sqrt{2E_0}$.
\end{enumerate}

\medskip
\noindent\textbf{Bagian VI --- Metodenya yang diekspor:
Lotka--Volterra.} Resep bandulnya --- yakni integral pertama,
kurva aras yang \omterm{def:b3:topology:compact}{kompak}, dan ketunggalan --- menyelesaikan sebuah ekosistem.
Tetapkan $\alpha, \beta, \gamma, \delta > 0$ lalu tinjaulah, pada
kuadran terbuka $Q = \intoo0{+\infty}\times\intoo0{+\infty}$,
\[
x' = x\,(\alpha - \beta y), \qquad
y' = y\,(\delta x - \gamma)
\]
(dengan $x$ mangsanya dan $y$ pemangsanya).
\begin{enumerate}[resume]
  \item Tunjukkan bahwa $Q$ bersifat invarian --- sebab sumbunya berupa gabungan
        \omterm{def:b3:groups:action}{orbit}, yang dapat dihitung secara eksplisit, yang tak boleh diseberangi
        solusi mana pun (\cref{cor:b3:ode:uniqueness}) --- dan
        bahwa satu-satunya kesetimbangan di $Q$ adalah $(x_*, y_*) =
        \bigl(\frac\gamma\delta, \frac\alpha\beta\bigr)$.
  \item Tunjukkan bahwa
        \[
        H(x, y) = \delta x - \gamma\ln x + \beta y -
        \alpha\ln y
        \]
        merupakan integral pertama, bahwa $H = f(x) + g(y)$ dengan
        $f, g$ yang cembung sejati dan meledak di tak hingga pada
        $\intoo0{+\infty}$ dengan minimum di $x_*$ dan $y_*$, lalu
        turunkan bahwa semua \omterm{thm:b3:ode:maximal}{solusi maksimal} di $Q$ bersifat global.
  \item Tunjukkan bahwa untuk $h > h_* = H(x_*, y_*)$ himpunan aras
        $\{H = h\}\cap Q$ merupakan kurva tertutup di sekitar
        kesetimbangannya: yakni dua cabang \omterm{def:b3:topology:continuity}{kontinu} $y_\pm(x)$ atas
        sebuah selang \omterm{def:b3:topology:compact}{kompak} $\intcc{x_-}{x_+} \ni x_*$, yang direkatkan
        di ujungnya --- yakni analogi oval bandulnya.
  \item Buktikan bahwa setiap \omterm{def:b3:groups:action}{orbit} takkesetimbangan di $Q$ bersifat
        periodik: tegakkanlah peredaran berlawanan arah jarum jam
        melalui keempat daerah yang teriris garis $x = x_*$
        dan $y = y_*$, lalu batasi waktu penyeberangan setiap busurnya
        oleh sebuah integral berkesingularan akar kuadrat yang
        konvergen (seperti pada pertanyaan 5), lalu tutuplah dengan
        ketunggalannya (seperti pada pertanyaan 4).
  \item (Hukum rata-rata Volterra) Jika $T$ periode
        sebuah \omterm{def:b3:groups:action}{orbit} semacam itu, tunjukkan bahwa
        \[
        \frac1T\int_0^Tx(t)\,\dd t = \frac\gamma\delta,
        \qquad
        \frac1T\int_0^Ty(t)\,\dd t = \frac\alpha\beta :
        \]
        yakni rata-rata waktunya sama dengan nilai kesetimbangannya,
        berapa pun amplitudonya \emph{(integralkanlah $(\ln x)' =
        \alpha - \beta y$ atas satu periode)}.
  \item (Paradoks penangkapan ikan) Panenlah kedua spesiesnya dengan laju
        $\varepsilon \in \intoo0\alpha$: maka sistemnya menjaga
        bentuknya dengan $\alpha - \varepsilon$ dan $\gamma +
        \varepsilon$ menggantikan $\alpha$ dan $\gamma$. Lalu apa
        yang terjadi pada populasi rata-ratanya? Terangkanlah
        pengamatan d'Ancona (1914--1918): bahwa ketika penangkapan ikan Adriatik
        menurun selama perang, proporsi pemangsanya (yakni hiu) pada
        tangkapannya justru \emph{naik} ---
        dan mengapa penangkapan yang sedang menguntungkan mangsanya.
  \item Tulislah moral sepuluh barisnya: teorema mana pada
        bab ini yang menggerakkan setiap langkahnya, apa yang menggantikan
        energi bandulnya, dan mengapa kedua sistemnya tak memerlukan ---
        atau tak memiliki --- solusi bentuk tertutup yang dasar.

  \item (Modulasi lajunya) Pada rezim rotasinya $E_0 > 1$,
        tunjukkan bahwa $y = x'$ berosilasi antara
        $\sqrt{2(E_0 - 1)}$ (di $x \equiv \pi \bmod 2\pi$)
        dan $\sqrt{2(E_0 + 1)}$ (di $x \equiv 0$), bahwa
        rata-rata waktu $y$ atas satu periode tepat
        $\frac{2\pi}{\tau(E_0)}$, dan bahwa nisbah
        modulasinya $\sqrt{\frac{E_0 + 1}{E_0 - 1}} \to 1$ saat $E_0
        \to \infty$: sehingga rotasi cepat bersifat asimtotik
        seragam.
  \item (Kemonotonan periode rotasinya) Tunjukkan bahwa
        $\tau(E_0)$ bersifat $\mathcal C^1$ dan turun
        sejati pada $\intoo1{+\infty}$
        \emph{(turunkanlah di bawah tanda integralnya, dengan
        dominasi pada setiap $\intco{1 + \delta}\infty$)},
        dengan $\tau \to \infty$ saat $E_0 \to 1^+$ dan $\tau
        \to 0$ saat $E_0 \to \infty$. Lalu rakitlah gambaran
        bifurkasi \omterm{def:b3:complete:complete}{lengkap} bandulnya sepanjang sumbu
        energinya: yakni kesetimbangan di $E_0 = -1$, librasi
        berperiode naik dari $2\pi$ ke $\infty$ pada
        $-1 < E_0 < 1$, separatriks di $E_0 = 1$, dan
        rotasi berperiode turun dari $\infty$ ke
        $0$ selebihnya.

\end{enumerate}
\end{problem}
